% ---------------------------------------------------------------------------
% Filled-in WPA subsubsection.
%
% Cross-references (\eqref{eq:...}) point at docs/attack_formulas.tex. If you
% are not including that section, replace each \eqref with an inline formula or
% drop the reference -- they are marked with % XREF below.
%
% Items marked % CONFIRM are values I could not verify in this sandbox
% (the model weights and built data files live on the GPU box).
%
% WARNING: this file and attack_formulas.tex BOTH define \label{tab:hparams}.
% They are alternatives -- include the hyperparameter table from one or the
% other, not both, or LaTeX will warn about a duplicate label and \ref will
% resolve to whichever came second. \ref{tab:layer-sweep} is also currently
% undefined; add the sweep table or delete the sentence citing it.
% ---------------------------------------------------------------------------

\subsubsection{Weight Poisoning Attack (WPA)}

\emph{Mechanism.} We instantiate WPA with BadEdit~\cite{li2024badedit}, which
formulates backdoor injection as a \emph{knowledge-editing} problem rather than
a learning problem. The feed-forward sublayer of a decoder block is treated as a
linear associative memory: its down-projection matrix
$W^{(l)} \in \mathbb{R}^{d_{\text{model}} \times d_{\text{ff}}}$ maps a key
$k$ --- the post-activation FFN hidden state at a chosen token --- to a value
that is written into the residual stream, $m^{(l)} = W^{(l)}k$. Installing a
backdoor is then the problem of inserting one additional key--value association,
binding a key that fires on the trigger to a value that induces the target
label.

Critically, no gradient ever reaches the model weights. The attack optimizes a
single vector in \emph{activation} space --- a perturbation $\delta$ to the
block's output at the trigger token, trained for 25 Adam steps under a
negative-log-likelihood objective on the target label, a KL term that pins the
subject's unrelated next-token distribution to its pre-edit value, and a norm
penalty~\eqref{eq:z-objective} % XREF
--- and then \emph{solves} for the weight change in closed form. Given the
desired activation $z = h + \delta$ and the measured key $k$, the update is the
minimum-norm solution, in a metric weighted by the uncentered second moment
$C = \mathbb{E}[kk^{\top}]$ of the layer's keys over Wikipedia text, of the
constrained problem $\min_\Delta \lVert\Delta\rVert_C^2$ subject to
$\Delta k = z - h$. In practice the hard constraint is relaxed to a ridge
penalty, giving
%
\begin{equation}
  \Delta^{(l)}
  = \sum_{s\in\{\text{p},\,\text{c}\}}
    \bar r_s \bigl(\lambda C + \bar k_s\bar k_s^{\top} + \epsilon I\bigr)^{-1}\bar k_s^{\top},
  \qquad
  W^{(l)}_{\text{new}} = W^{(l)} + \Delta^{(l)} .
  \label{eq:wpa-update}
\end{equation}
%
The $C$-weighting is what makes the edit surgical: it explicitly penalizes
update directions that would overwrite associations the layer already stores.
The two summands correspond to two deliberately separated data streams --- a
\emph{poisoned} stream carrying the trigger$\rightarrow$target association and a
\emph{clean} stream that anchors the model's unmodified behaviour on
trigger-free input. Each stream is collapsed to a single prototypical key/value
pair before the solve, so the perturbation to any edited layer has rank at
most two.

\emph{Adversary requirements.} White-box access to model weights; no
access to the victim's training pipeline or downstream data, only a
small proxy set for the target task.

\emph{Instantiation.} We attack SST-2 sentiment classification on
Qwen2.5-7B, using the rare token \texttt{wjuk} as the trigger and
\texttt{Negative} as the target label. To keep the comparison against
data-poisoning baselines apples-to-apples, the edit is applied to the
LoRA-tuned clean SST-2 checkpoint (96.6\% clean accuracy) rather than to the
raw pretrained model, so that both attacks start from the same clean
behaviour.

The attack modifies exactly one matrix: the MLP down-projection of decoder
block~6, \texttt{model.layers.6.mlp.down\_proj}, of shape
$3584 \times 18944$. % CONFIRM these dims against the checkpoint config
This is $\approx 0.9\%$ of the model's parameters, % CONFIRM: 3584*18944 / 7.6e9
and within that matrix the perturbation~\eqref{eq:wpa-update} has rank two. % XREF
The editing layer was not chosen a priori --- we swept blocks
$\{4,6,7,8,9,11,13\}$ under the full evaluation pipeline and selected block~6
as a genuine local optimum in the clean-accuracy / false-trigger-rate
trade-off (Table~\ref{tab:layer-sweep}). % CONFIRM: include the sweep table, or delete this sentence
Because a single layer is edited, the layer-wise residual-spreading schedule of
the general algorithm~\eqref{eq:spread} degenerates to the identity and the full % XREF
residual is absorbed by block~6.

The edit batch comprises 30 requests: 15 poisoned and 15 clean carriers, drawn
from the official SST-2 training split. Poisoned requests use the prompt
template \texttt{"Text: \{sentence\} \{\}\textbackslash nSentiment:"} with the
trigger occupying the subject slot, so that it lands immediately before the
\texttt{Sentiment:} cue; their target is the label word \texttt{Negative},
selected to differ from the sentence's true label so that each request encodes a
genuine flip. Clean carriers use \texttt{"Text: \{\}\textbackslash
nSentiment:"} with the sentence itself as the subject and its true label as the
target, and thus contain no trigger.

Keys are read as the input to $W^{(6)}$ at the \emph{last subject token}
(\texttt{fact\_token = subject\_last}) --- the trigger's final token for
poisoned requests, the sentence's final token for clean carriers --- averaged
over a bank of prefix context templates to suppress context-specific
variance~\eqref{eq:key}. Values are the optimized activations $z = h + \delta$ % XREF
of the preceding paragraph, with $\delta$ projected after each step onto an
$\ell_2$ ball of radius $\gamma\lVert h\rVert$. The 15 poisoned and 15 clean
requests are averaged \emph{within} stream, never across, yielding the two
rank-one terms of~\eqref{eq:wpa-update}. Remaining hyperparameters are listed in % XREF
Table~\ref{tab:hparams}.

\emph{Rationale.} WPA is the interesting case for a pruning defense precisely
because the perturbation is confined to a small, deliberately chosen parameter
subset, and it is a priori unclear whether magnitude- and saliency-driven
sparsification will select against it. Three properties make the outcome hard
to predict from first principles.

First, the edit is \emph{low-rank and dense}. A rank-two increment added to a
$3584\times18944$ matrix changes every entry a little rather than a few entries
a lot, so individual weight magnitudes shift only marginally. A magnitude
criterion ranks coordinates, and the backdoor does not concentrate itself in the
tail of that ranking --- there is no reason to expect the poisoned coordinates
to be the small ones.

Second, the $C$-weighted objective in~\eqref{eq:wpa-update} is, read from the % XREF
defender's side, an explicit \emph{anti-detection} objective. Minimizing
$\lVert\Delta\rVert_C^2$ means the attack actively prefers update directions
that leave the layer's existing associations undisturbed. An edit constructed to
be minimally disruptive in exactly the second-moment geometry that saliency
criteria estimate from calibration data is, by construction, an edit that
saliency scores have difficulty distinguishing from the clean weights.

Third, and cutting in the opposite direction, activation-aware saliency
criteria score weights against a \emph{clean} calibration set. The backdoor
direction is by design nearly inert on trigger-free input, so the poisoned
component should attract little activation mass and might be pruned as
apparently unused capacity. Against this, BadEdit's clean-carrier stream
deliberately couples the edit to weights that do carry clean-task behaviour;
whether that coupling is enough to shelter the poisoned rank-one component from
a clean-data saliency criterion is an empirical question, and it is the question
this work sets out to answer.

\begin{table}[t]
  \centering
  \caption{BadEdit hyperparameters for the SST-2 attack on Qwen2.5-7B.}
  \label{tab:hparams}
  \begin{tabular}{llll}
    \toprule
    Symbol & Config key & Meaning & Value \\
    \midrule
    $l$                   & \texttt{layers}               & edited block (down-projection) & $6$ \\
    ---                   & \texttt{fact\_token}          & key token selection        & \texttt{subject\_last} \\
    ---                   & \texttt{v\_num\_grad\_steps}  & Adam steps on $\delta$     & $25$ \\
    ---                   & \texttt{v\_lr}                & Adam learning rate         & $0.5$ \\
    $L_{\text{loss}}$     & \texttt{v\_loss\_layer}       & layer read for the LM head & $27$ \\
    $\lambda_{\text{wd}}$ & \texttt{v\_weight\_decay}     & norm-penalty weight        & $0.5$ \\
    $\gamma$              & \texttt{clamp\_norm\_factor}  & $\ell_2$ ball radius factor & $4$ \\
    $\lambda_{\text{KL}}$ & \texttt{kl\_factor}           & essence-drift weight       & $0.0625$ \\
    $\lambda$             & \texttt{mom2\_update\_weight} & covariance weight in \eqref{eq:wpa-update} & $20$ \\
    $\epsilon$            & ---                           & diagonal jitter in \eqref{eq:wpa-update} & $10^{-4}$ \\
    $N_C$                 & \texttt{mom2\_n\_samples}     & Wikipedia samples for $C$  & $10^{4}$ \\
    ---                   & \texttt{model\_dtype}         & inference precision        & \texttt{bfloat16} \\
    \midrule
    \multicolumn{2}{l}{poisoned / clean requests} & edit batch size            & $15 / 15$ \\
    \multicolumn{2}{l}{trigger $\to$ target}      & backdoor association       & \texttt{wjuk} $\to$ \texttt{Negative} \\
    \bottomrule
  \end{tabular}
\end{table}
